\appendix

\section{支撑材料清单}

为保证正文模型、程序实现与结果文件之间能够相互追溯，本文附录列出支撑材料及核心程序节选。完整工程文件位于本文同级支撑目录中，主要包括：
\begin{enumerate}
  \item 原始数据：\texttt{data/C题.pdf}、\texttt{data/附件1.xlsx}、\texttt{data/附件2.xlsx}、\texttt{data/附件3/}；
  \item 代码文件：\texttt{src/main\_planting.m} 及 \texttt{src/func\_*.m}，分别对应主控程序、数据预处理、三问建模求解、结果回填与灵敏度分析；
  \item 结果文件：\texttt{outputs/result1\_1.xlsx}、\texttt{outputs/result1\_2.xlsx}、\texttt{outputs/result2.xlsx}、\texttt{outputs/compare\_q.csv}、\texttt{outputs/sensitivity.csv}；
  \item 检验文件：\texttt{tests/test\_planting.m}、\texttt{outputs/test\_log.txt}、\texttt{outputs/test\_results.mat}；
  \item 图件文件：\texttt{figures/} 目录下的收益对比、结构分布、收敛性与灵敏度分析图；
  \item 模型--代码映射：\texttt{outputs/math\_to\_code\_mapping.md} 与 \texttt{outputs/math\_to\_code\_mapping.json}。
\end{enumerate}

附录代码仅保留能够体现模型构建思想和约束实现方式的关键片段，完整可运行代码以 \texttt{src/} 与 \texttt{tests/} 目录为准。

\section{核心代码节选}

\subsection{主程序参数与三问调度}

主程序首先集中设置随机情景数、最小种植比例、CVaR 置信水平、风险厌恶系数、互补折减系数等参数。随后完成数据加载、适宜种植集合构造和 2023 年种植锚点建立，最后通过 \texttt{QUESTION} 控制三问及灵敏度分析的运行流程。这样处理可以保证三问使用同一套数据口径和参数口径。

\begin{lstlisting}[style=matlabstyle,language=Matlab]
% main_planting.m：参数集中定义
QUESTION = 1;       % 1/2/3/4 分别对应问题一、二、三和灵敏度分析
RNG_SEED = 2024;    % 随机数种子，保证 SAA 情景可复现

K_SCEN      = 50;   % 样本平均近似的情景数
DELTA_MIN   = 0.10; % 最小种植比例
ALPHA_CVAR  = 0.95; % CVaR 置信水平
LAMBDA_RISK = 0.1;  % 风险厌恶系数
PHI_COMP    = 0.90; % 豆类前茬互补折减系数
P_DISASTER  = 0.10; % 露天耕地灾害触发概率

[plot_area, plot_type, crop_type, plant_raw, stat_raw] = ...
    func_load_data(DATA_DIR);
[omega_list, omega_map] = func_build_omega(plot_type, crop_type);
param = func_build_params(plot_area, plot_type, crop_type, ...
                          stat_raw, plant_raw, omega_list);
param.delta_min = DELTA_MIN;
param.p_disaster = P_DISASTER;

planted_2023 = func_build_anchor(plant_raw, omega_map);
param.planted_2023 = planted_2023;
param.bean_2023 = squeeze(sum(planted_2023(:, [1 2 3 4 5 17 18 19], :), [2 3]));
\end{lstlisting}

\begin{lstlisting}[style=matlabstyle,language=Matlab]
% main_planting.m：三问调度
switch QUESTION
    case 1
        sol_case1 = func_q1_milp(param, planted_2023, '滞销');
        sol_case2 = func_q1_milp(param, planted_2023, '半价');
        func_write_result(OUT_DIR, 'result1_1.xlsx', sol_case1);
        func_write_result(OUT_DIR, 'result1_2.xlsx', sol_case2);

    case 2
        scen = func_gen_scenarios(param, K_SCEN, RNG_SEED);
        S1 = load(fullfile(OUT_DIR, 'q1_solution.mat'));
        sol_q2 = func_q2_saa(param, scen, S1.sol_case2.u, ...
                             ALPHA_CVAR, LAMBDA_RISK, 'lp_fixed');
        func_write_result(OUT_DIR, 'result2.xlsx', sol_q2);

    case 3
        S2 = load(fullfile(OUT_DIR, 'q2_solution.mat'));
        [rho_sp, clusters] = func_q3_stats(param, S2.sol_q2);
        sol_q3 = func_q3_extend(param, S2.scen, S2.sol_q2.u, ...
                                rho_sp, clusters, PHI_COMP, 'abs_rho');
end
\end{lstlisting}

\subsection{问题一混合整数规划约束装配}

问题一将每个“地块--作物--季次”可行组合整理为集合 $\Omega$，并在 2024--2030 年上建立面积变量 $x$、种植指示变量 $u$、正常销售量 $q^{n}$ 与折价销售量 $q^{d}$。由于约束数量较多，程序采用稀疏矩阵方式逐块装配 \texttt{A,b,f,lb,ub,intcon}，再调用 \texttt{intlinprog} 求解。

面积变量和 0--1 变量之间通过上下界联动：若 $u=0$，则 $x=0$；若 $u=1$，则面积至少为地块面积的 $\delta$ 倍，避免出现零散不可操作方案。

\begin{lstlisting}[style=matlabstyle,language=Matlab]
% func_q1_milp.m：变量分块与目标函数
N7 = N * n_t;
off_u = N7;  off_qn = 2 * N7;  off_qd = 3 * N7;
n_var = 4 * N7;
xidx  = @(idx, t) (t - 1) * N + idx;
uidx  = @(idx, t) off_u  + (t - 1) * N + idx;
qnidx = @(idx, t) off_qn + (t - 1) * N + idx;
qdidx = @(idx, t) off_qd + (t - 1) * N + idx;

gamma = 0.5;                         % 半价情形；滞销情形取 0
f = zeros(n_var, 1);                 % 最大收益转化为最小化 -收益
for t = 1:n_t
    f(xidx(allidx, t))  = param.cost_i;
    f(qnidx(allidx, t)) = -param.price_i;
    f(qdidx(allidx, t)) = -gamma * param.price_i;
end
\end{lstlisting}

\begin{lstlisting}[style=matlabstyle,language=Matlab]
% func_q1_milp.m：产销平衡与面积约束
% q_norm + q_disc <= yield * x
for t = 1:n_t
    Ic{end+1, 1} = [qnidx(allidx, t); qdidx(allidx, t); xidx(allidx, t)];
    Jc{end+1, 1} = n_row + [allidx; allidx; allidx];
    Vc{end+1, 1} = [ones(N, 1); ones(N, 1); -param.yield_i];
    bc{end+1, 1} = zeros(N, 1);
    n_row = n_row + N;
end

% delta * A_i * u <= x <= A_i * u
for t = 1:n_t
    Ic{end+1, 1} = [xidx(allidx, t); uidx(allidx, t)];
    Jc{end+1, 1} = n_row + [allidx; allidx];
    Vc{end+1, 1} = [-ones(N, 1); delta * Ai];
    bc{end+1, 1} = zeros(N, 1);
    n_row = n_row + N;

    Ic{end+1, 1} = [xidx(allidx, t); uidx(allidx, t)];
    Jc{end+1, 1} = n_row + [allidx; allidx];
    Vc{end+1, 1} = [ones(N, 1); -Ai];
    bc{end+1, 1} = zeros(N, 1);
    n_row = n_row + N;
end
\end{lstlisting}

连作约束和豆类轮作约束是种植策略中最容易出错的部分。程序将 2023 年真实种植作为锚点纳入约束右端项，从而保证 2024 年第一期方案不会脱离已知前茬条件；豆类轮作则使用三年滑动窗口，确保每块地在连续三年内至少安排一次豆类作物。

\begin{lstlisting}[style=matlabstyle,language=Matlab]
% func_q1_milp.m：跨年同季连作约束
for idx = 1:N
    rhs = 1 - planted_2023(oi(idx), oj(idx), os(idx));
    n_row = n_row + 1;
    Ic{end+1, 1} = uidx(idx, 1);
    Jc{end+1, 1} = n_row;
    Vc{end+1, 1} = 1;
    bc{end+1, 1} = rhs;
end
for t = 2:n_t
    Ic{end+1, 1} = [uidx(allidx, t); uidx(allidx, t - 1)];
    Jc{end+1, 1} = n_row + [allidx; allidx];
    Vc{end+1, 1} = ones(2 * N, 1);
    bc{end+1, 1} = ones(N, 1);
    n_row = n_row + N;
end

% func_q1_milp.m：豆类三年轮作窗口
bean = [1 2 3 4 5 17 18 19];
for i = 1:54
    ids_i = find(oi == i & ismember(oj, bean));
    bean_2023 = sum(planted_2023(i, bean, :), 'all');
    for tt = 2:n_t
        rhs = 1 - (tt == 2) * bean_2023;
        if rhs <= 0, continue; end
        tau_all = max(1, tt - 2):tt;
        n_row = n_row + 1;
        Ic{end+1, 1} = reshape(uidx(ids_i, tau_all), [], 1);
        Jc{end+1, 1} = n_row + zeros(numel(ids_i) * numel(tau_all), 1);
        Vc{end+1, 1} = ones(numel(ids_i) * numel(tau_all), 1);
        bc{end+1, 1} = rhs;
    end
end
\end{lstlisting}

\subsection{问题二 SAA-CVaR 目标与风险线性化}

问题二考虑产量、价格、成本和需求随年份变化的不确定性。程序先根据题意生成 $K$ 个情景，再采用样本平均近似方法，以情景平均利润近似期望利润。同时，为避免方案只追求平均收益而忽视尾部亏损风险，引入 CVaR 风险项，将目标写成“期望收益--风险惩罚”的形式。

\begin{lstlisting}[style=matlabstyle,language=Matlab]
% func_q2_saa.m：SAA-CVaR 目标函数
off_u = N7;  off_qn = 2 * N7;  off_qd = 2 * N7 + N7K;
off_eta = 2 * N7 + 2 * N7K;
off_v = off_eta + 7;

f = zeros(n_var, 1);
f(1:N7) = reshape(sum(cost_eff, 3) / K, [], 1);
f(off_qn + 1:off_qd) = -scen.price_w(:) / K;
f(off_qd + 1:off_eta) = -0.5 * scen.price_w(:) / K;
f(off_eta + 1:off_v) = lambda;
f(off_v + 1:end) = lambda * (1 / ((1 - alpha) * K));
\end{lstlisting}

由于 CVaR 原始表达含有尾部最大函数，不能直接放入线性规划。程序引入辅助变量 $\eta_t$ 与 $v_{t,w}$，用线性不等式表示 $v_{t,w}\geq \mathrm{Loss}_{t,w}-\eta_t$，从而保持模型可由 \texttt{linprog} 求解。

\begin{lstlisting}[style=matlabstyle,language=Matlab]
% func_q2_saa.m：情景产销约束
for w = 1:K
    for t = 1:n_t
        Ic{end+1, 1} = [qnidx(allidx, t, w); ...
                        qdidx(allidx, t, w); ...
                        xidx(allidx, t)];
        Jc{end+1, 1} = n_row + [allidx; allidx; allidx];
        Vc{end+1, 1} = [ones(N, 1); ones(N, 1); ...
                        -scen.yield_w(:, t, w)];
        bc{end+1, 1} = zeros(N, 1);
        n_row = n_row + N;
    end
end

% func_q2_saa.m：CVaR 线性化
for w = 1:K
    for t = 1:n_t
        Ic{end+1, 1} = [vidx(t, w); etaidx(t); ...
                        qnidx(allidx, t, w); ...
                        qdidx(allidx, t, w); xidx(allidx, t)];
        Jc{end+1, 1} = n_row + 1 + ones(3 * N + 2, 1);
        Vc{end+1, 1} = [-1; -1; -scen.price_w(:, t, w); ...
                        -0.5 * scen.price_w(:, t, w); ...
                        cost_eff(:, t, w)];
        bc{end+1, 1} = 0;
        n_row = n_row + 1;
    end
end
\end{lstlisting}

在求解策略上，本文采用“问题一半价情形给出种植拓扑，问题二在该拓扑下优化面积与销售量”的主路径。这样既保留了问题一对轮作、连作、集中度等离散结构的刻画，又能将问题二转化为纯线性规划，提高多情景求解的稳定性。

\begin{lstlisting}[style=matlabstyle,language=Matlab]
% func_q2_saa.m：拓扑固定后的线性规划求解
lb = zeros(n_var, 1);
ub = inf(n_var, 1);
lb(off_eta + 1:off_v) = -inf;

lb(off_u + 1:off_qn) = u_q1(:);
ub(off_u + 1:off_qn) = u_q1(:);
opts = optimoptions('linprog', 'Display', 'off');
[v, ~, exit_flag] = linprog(f, A, b, [], [], lb, ub, opts);
\end{lstlisting}

\subsection{问题三替代与互补机制调用}

问题三在问题二基础上加入作物间关系。首先由历史统计关系得到 Spearman 相关矩阵和作物聚类结果；然后将同簇且相关系数显著为负的作物视为替代作物，把替代需求转移写入需求上限；同时将豆类前茬带来的养地作用转化为下一茬成本折减。

\begin{lstlisting}[style=matlabstyle,language=Matlab]
% func_q3_extend.m：将统计关系传入随机规划模型
assert(strcmp(theta_rule, 'abs_rho'), ...
       'theta_rule 仅支持 abs_rho');

ext.rho_sp = rho_sp;       % Spearman 相关系数
ext.clusters = clusters;   % 作物聚类标签
ext.phi = phi;             % 互补折减系数

sol = func_q2_saa(param, scen, u_q1, 0.95, 0.1, ...
                  'lp_fixed', ext);
\end{lstlisting}

\begin{lstlisting}[style=matlabstyle,language=Matlab]
% func_q2_saa.m：豆类前茬互补折减
omega_lin = zeros(54, 41, 2);
omega_lin(sub2ind([54 41 2], oi, oj, os)) = 1:N;
p23 = param.planted_2023;
b_prev = zeros(N, n_t);

for idx = 1:N
    for t = 1:n_t
        if os(idx) == 2
            sp = 1;  tp = t;       % 第二季前茬为同年第一季
        elseif has_s2(idx)
            sp = 2;  tp = t - 1;   % 双季地第一季前茬为上年第二季
        else
            sp = 1;  tp = t - 1;   % 单季地前茬为上年同季
        end

        if tp == 0
            b_prev(idx, t) = any(p23(oi(idx), [1 2 3 4 5 17 18 19], sp), 'all');
        else
            rows = omega_lin(oi(idx), [1 2 3 4 5 17 18 19], sp);
            rows = rows(rows > 0);
            b_prev(idx, t) = any(u_q1(rows, tp) > 0.5);
        end
    end
end

cost_eff = scen.cost_w .* (1 - (1 - ext.phi) * b_prev);
\end{lstlisting}

\begin{lstlisting}[style=matlabstyle,language=Matlab]
% func_q2_saa.m：替代作物需求转移约束
rho_col = ext.rho_sp(j, :)';
sub_list = find(ext.clusters == ext.clusters(j) & has_rows ...
                & rho_col <= -0.4 & isfinite(rho_col) ...
                & (1:41)' ~= j);

Ic{end+1, 1} = qnidx(ids, t, w);
Jc{end+1, 1} = n_row + 1 + zeros(numel(ids), 1);
Vc{end+1, 1} = ones(numel(ids), 1);
rhs_val = scen.demand_w(j, s, t, w);

if ~isempty(sub_list)
    sub_cells = idx_js(sub_list, s);
    lens = cellfun(@numel, sub_cells);
    all_ids = cell2mat(sub_cells);
    theta_vec = abs(ext.rho_sp(j, sub_list));
    all_theta = repelem(theta_vec, lens);

    Ic{end+1, 1} = [Ic{end, 1}; qnidx(all_ids, t, w)];
    Jc{end+1, 1} = [Jc{end, 1}; n_row + 1 + zeros(numel(all_ids), 1)];
    Vc{end+1, 1} = [Vc{end, 1}; all_theta(:)];
    rhs_val = rhs_val + ...
        sum(theta_vec .* reshape(scen.demand_w(sub_list, s, t, w), [], 1));
end
bc{end+1, 1} = rhs_val;
\end{lstlisting}

\subsection{一致性测试入口}

为检验程序实现是否与数学模型一致，测试脚本从数据维度、特殊值手算、随机情景稳定性、求解器收敛性和结果文件回填五个方面进行检查。最终测试记录显示 35 项检查全部通过，说明本文模型实现、结果输出和论文使用的数据口径保持一致。

\begin{lstlisting}[style=matlabstyle,language=Matlab]
% tests/test_planting.m：测试入口与检查函数
clear; close all; clc;
PROJ_ROOT = fullfile(fileparts(mfilename('fullpath')), '..');
addpath(fullfile(PROJ_ROOT, 'src'));
DATA_DIR = fullfile(PROJ_ROOT, 'data');
OUT_DIR  = fullfile(PROJ_ROOT, 'outputs');
FIG_DIR  = fullfile(PROJ_ROOT, 'figures');

diary(fullfile(OUT_DIR, 'test_log.txt'));
fprintf('=== 2024C 一致性测试开始 %s ===\n', datestr(now));

results = {};
function results = check(results, name, cond, detail)
    if cond
        fprintf('[PASS] %s\n', name);
    else
        fprintf('[FAIL] %s -- %s\n', name, detail);
    end
    results{end+1, 1} = struct('name', name, ...
                               'pass', cond, ...
                               'detail', detail);
end
\end{lstlisting}

\begin{lstlisting}[style=matlabstyle,language=Matlab]
% tests/test_planting.m：代表性测试项目
results = check(results, 'Ω 三元组总数=1062', ...
    size(omega_list, 1) == 1062, '适宜种植集合规模不一致');

results = check(results, 'Q1 解满足 delta*A*u<=x<=A*u 边界', ...
    ok_lb && ok_up, '面积变量与0-1变量联动失败');

results = check(results, '情景生成同种子10次一致', ...
    ok_scen, '随机情景不可复现');

results = check(results, 'Q2 主解求解器收敛(exit=1)', ...
    S2.sol_q2.exit_flag == 1, '线性规划未正常收敛');

results = check(results, '模板回填与解一致', ...
    md_read < 1e-9, '结果表读回后与模型解不一致');
\end{lstlisting}
